\documentclass[12pt,letterpaper]{article}
\usepackage[margin=0.65in]{geometry}
\usepackage[T1]{fontenc}
\usepackage{lmodern}
\pdfmapfile{+lm.map}
\usepackage{tikz}
\pagestyle{empty}
\setlength{\parindent}{0pt}
\hyphenpenalty=10000\exhyphenpenalty=10000
\begin{document}
\begin{tikzpicture}[remember picture,overlay]\begin{scope}[shift={(current page.north west)},x=1in,y=-1in]
\node[anchor=north west,inner sep=0pt,text width=7.2in,align=left,font=\fontsize{11}{13.75}\selectfont] at (0.65,0.41) {Week 19 / No card inside another / K--1};
\draw[line width=.35pt] (.65,.70)--(7.85,.70);
\node[anchor=north west,inner sep=0pt,text width=6.7in,align=left,font=\fontsize{9}{11.25}\selectfont] at (0.65,10.43) {Bellingham Math Circle / Week 19 / F19-K-v2};
\node[anchor=north west,inner sep=0pt,text width=0.55in,align=right,font=\fontsize{9}{11.25}\selectfont] at (7.3,10.43) {1};
\node[anchor=north west,inner sep=0pt,text width=7.2in,align=left,font=\fontsize{15}{18.75}\selectfont] at (0.65,0.96) {A card fits inside another when all its pictures are on the other card. The empty card fits inside every card. In a collection, no card may fit inside another. Use a card only once in a collection.};
\node[anchor=north west,inner sep=0pt,text width=7.2in,align=left,font=\fontsize{16}{20.0}\selectfont] at (0.65,2.25) {\textbf{Problem 1:} In each group, circle as many cards as you can without breaking the rule.};
\draw[line width=.7pt,rounded corners=2pt] (2.41,3.0) rectangle (3.35,3.94);
\draw[line width=1.5pt] (2.4852000000000003,3.0) -- (2.5604,3.0);
\draw[line width=.7pt,rounded corners=2pt] (3.7800000000000002,3.0) rectangle (4.720000000000001,3.94);
\draw[line width=1.5pt] (3.8552000000000004,3.0) -- (3.9304,3.0);
\draw[line width=.95pt] (4.043200000000001,3.2632) circle[radius=0.11467999999999999in];
\draw[line width=.7pt,rounded corners=2pt] (5.15,3.0) rectangle (6.09,3.94);
\draw[line width=1.5pt] (5.2252,3.0) -- (5.300400000000001,3.0);
\draw[line width=.95pt,line join=round] (5.8268,3.131318)--(5.952948,3.360678)--(5.700652000000001,3.360678)--cycle;
\draw[line width=.7pt,rounded corners=2pt] (2.41,4.55) rectangle (3.35,5.49);
\draw[line width=1.5pt] (2.4852000000000003,4.55) -- (2.5604,4.55);
\draw[line width=.95pt] (2.6732,4.8132) circle[radius=0.11467999999999999in];
\draw[line width=.7pt,rounded corners=2pt] (3.7800000000000002,4.55) rectangle (4.720000000000001,5.49);
\draw[line width=1.5pt] (3.8552000000000004,4.55) -- (3.9304,4.55);
\draw[line width=.95pt] (4.043200000000001,4.8132) circle[radius=0.11467999999999999in];
\draw[line width=.95pt,line join=round] (4.4568,4.681318)--(4.582948,4.910678)--(4.330652000000001,4.910678)--cycle;
\draw[line width=.7pt,rounded corners=2pt] (5.15,4.55) rectangle (6.09,5.49);
\draw[line width=1.5pt] (5.2252,4.55) -- (5.300400000000001,4.55);
\draw[line width=.95pt,line join=round] (5.8268,4.681318)--(5.952948,4.910678)--(5.700652000000001,4.910678)--cycle;
\draw[line width=.7pt,rounded corners=2pt] (2.41,6.1) rectangle (3.35,7.039999999999999);
\draw[line width=1.5pt] (2.4852000000000003,6.1) -- (2.5604,6.1);
\draw[line width=.95pt] (2.6732,6.3632) circle[radius=0.11467999999999999in];
\draw[line width=.95pt,line join=round] (3.0868,6.231318)--(3.2129480000000004,6.460678)--(2.960652,6.460678)--cycle;
\draw[line width=.7pt,rounded corners=2pt] (3.7800000000000002,6.1) rectangle (4.720000000000001,7.039999999999999);
\draw[line width=1.5pt] (3.8552000000000004,6.1) -- (3.9304,6.1);
\draw[line width=.95pt] (4.043200000000001,6.3632) circle[radius=0.11467999999999999in];
\draw[line width=.95pt] (3.9285200000000007,6.66212) rectangle (4.1578800000000005,6.89148);
\draw[line width=.7pt,rounded corners=2pt] (5.15,6.1) rectangle (6.09,7.039999999999999);
\draw[line width=1.5pt] (5.2252,6.1) -- (5.300400000000001,6.1);
\draw[line width=.95pt,line join=round] (5.8268,6.231318)--(5.952948,6.460678)--(5.700652000000001,6.460678)--cycle;
\draw[line width=.95pt] (5.298520000000001,6.66212) rectangle (5.527880000000001,6.89148);
\draw[line width=.7pt,rounded corners=2pt] (2.41,7.65) rectangle (3.35,8.59);
\draw[line width=1.5pt] (2.4852000000000003,7.65) -- (2.5604,7.65);
\draw[line width=.95pt] (2.6732,7.913200000000001) circle[radius=0.11467999999999999in];
\draw[line width=.7pt,rounded corners=2pt] (3.7800000000000002,7.65) rectangle (4.720000000000001,8.59);
\draw[line width=1.5pt] (3.8552000000000004,7.65) -- (3.9304,7.65);
\draw[line width=.95pt,line join=round] (4.4568,7.781318000000001)--(4.582948,8.010678)--(4.330652000000001,8.010678)--cycle;
\draw[line width=.95pt] (3.9285200000000007,8.21212) rectangle (4.1578800000000005,8.44148);
\draw[line width=.7pt,rounded corners=2pt] (5.15,7.65) rectangle (6.09,8.59);
\draw[line width=1.5pt] (5.2252,7.65) -- (5.300400000000001,7.65);
\draw[line width=.95pt] (5.413200000000001,7.913200000000001) circle[radius=0.11467999999999999in];
\draw[line width=.95pt,line join=round] (5.8268,7.781318000000001)--(5.952948,8.010678)--(5.700652000000001,8.010678)--cycle;
\draw[line width=.95pt] (5.298520000000001,8.21212) rectangle (5.527880000000001,8.44148);
\end{scope}\end{tikzpicture}\mbox{}
\newpage
\begin{tikzpicture}[remember picture,overlay]\begin{scope}[shift={(current page.north west)},x=1in,y=-1in]
\node[anchor=north west,inner sep=0pt,text width=7.2in,align=left,font=\fontsize{11}{13.75}\selectfont] at (0.65,0.41) {Week 19 / No card inside another / K--1};
\draw[line width=.35pt] (.65,.70)--(7.85,.70);
\node[anchor=north west,inner sep=0pt,text width=6.7in,align=left,font=\fontsize{9}{11.25}\selectfont] at (0.65,10.43) {Bellingham Math Circle / Week 19 / F19-K-v2};
\node[anchor=north west,inner sep=0pt,text width=0.55in,align=right,font=\fontsize{9}{11.25}\selectfont] at (7.3,10.43) {2};
\node[anchor=north west,inner sep=0pt,text width=7.2in,align=left,font=\fontsize{16}{20.0}\selectfont] at (0.65,0.98) {\textbf{Problem 2:} Find every collection of one or more cards from this deck.};
\draw[line width=.7pt,rounded corners=2pt] (2.17,1.79) rectangle (3.0,2.62);
\draw[line width=1.5pt] (2.2363999999999997,1.79) -- (2.3028,1.79);
\draw[line width=.7pt,rounded corners=2pt] (3.28,1.79) rectangle (4.109999999999999,2.62);
\draw[line width=1.5pt] (3.3463999999999996,1.79) -- (3.4128,1.79);
\draw[line width=.95pt] (3.5124,2.0224) circle[radius=0.10125999999999999in];
\draw[line width=.95pt,line join=round] (3.8775999999999997,1.9059510000000002)--(3.9889859999999997,2.108471)--(3.7662139999999997,2.108471)--cycle;
\draw[line width=.7pt,rounded corners=2pt] (4.39,1.79) rectangle (5.22,2.62);
\draw[line width=1.5pt] (4.4563999999999995,1.79) -- (4.522799999999999,1.79);
\draw[line width=.95pt] (4.6224,2.0224) circle[radius=0.10125999999999999in];
\draw[line width=.7pt,rounded corners=2pt] (5.5,1.79) rectangle (6.33,2.62);
\draw[line width=1.5pt] (5.5664,1.79) -- (5.6328,1.79);
\draw[line width=.95pt,line join=round] (6.0976,1.9059510000000002)--(6.208986,2.108471)--(5.9862139999999995,2.108471)--cycle;
\draw[black!22,line width=.35pt] (.65,5.12)--(7.85,5.12);
\node[anchor=north west,inner sep=0pt,text width=7.2in,align=left,font=\fontsize{16}{20.0}\selectfont] at (0.65,5.42) {\textbf{Problem 3:} Find as many different two-card collections as you can from this deck.};
\draw[line width=.7pt,rounded corners=2pt] (0.79,6.26) rectangle (1.48,6.949999999999999);
\draw[line width=1.5pt] (0.8452000000000001,6.26) -- (0.9004000000000001,6.26);
\draw[line width=.95pt] (0.9832000000000001,6.4532) circle[radius=0.08417999999999999in];
\draw[line width=.95pt] (0.89902,6.67262) rectangle (1.06738,6.84098);
\draw[line width=.7pt,rounded corners=2pt] (1.68,6.26) rectangle (2.37,6.949999999999999);
\draw[line width=1.5pt] (1.7351999999999999,6.26) -- (1.7904,6.26);
\draw[line width=.7pt,rounded corners=2pt] (2.57,6.26) rectangle (3.26,6.949999999999999);
\draw[line width=1.5pt] (2.6252,6.26) -- (2.6803999999999997,6.26);
\draw[line width=.95pt,line join=round] (3.0667999999999997,6.356393)--(3.159398,6.524753)--(2.9742019999999996,6.524753)--cycle;
\draw[line width=.7pt,rounded corners=2pt] (3.46,6.26) rectangle (4.15,6.949999999999999);
\draw[line width=1.5pt] (3.5152,6.26) -- (3.5704,6.26);
\draw[line width=.95pt] (3.6532,6.4532) circle[radius=0.08417999999999999in];
\draw[line width=.95pt,line join=round] (3.9568,6.356393)--(4.049398,6.524753)--(3.8642019999999997,6.524753)--cycle;
\draw[line width=.95pt] (3.56902,6.67262) rectangle (3.73738,6.84098);
\draw[line width=.7pt,rounded corners=2pt] (4.35,6.26) rectangle (5.039999999999999,6.949999999999999);
\draw[line width=1.5pt] (4.4052,6.26) -- (4.4604,6.26);
\draw[line width=.95pt] (4.5432,6.4532) circle[radius=0.08417999999999999in];
\draw[line width=.7pt,rounded corners=2pt] (5.239999999999999,6.26) rectangle (5.93,6.949999999999999);
\draw[line width=1.5pt] (5.2951999999999995,6.26) -- (5.3504,6.26);
\draw[line width=.95pt,line join=round] (5.7368,6.356393)--(5.829397999999999,6.524753)--(5.644202,6.524753)--cycle;
\draw[line width=.95pt] (5.349019999999999,6.67262) rectangle (5.517379999999999,6.84098);
\draw[line width=.7pt,rounded corners=2pt] (6.13,6.26) rectangle (6.82,6.949999999999999);
\draw[line width=1.5pt] (6.1852,6.26) -- (6.2404,6.26);
\draw[line width=.95pt] (6.3232,6.4532) circle[radius=0.08417999999999999in];
\draw[line width=.95pt,line join=round] (6.6268,6.356393)--(6.719398,6.524753)--(6.5342020000000005,6.524753)--cycle;
\draw[line width=.7pt,rounded corners=2pt] (7.02,6.26) rectangle (7.709999999999999,6.949999999999999);
\draw[line width=1.5pt] (7.0752,6.26) -- (7.1304,6.26);
\draw[line width=.95pt] (7.12902,6.67262) rectangle (7.2973799999999995,6.84098);
\end{scope}\end{tikzpicture}\mbox{}
\newpage
\begin{tikzpicture}[remember picture,overlay]\begin{scope}[shift={(current page.north west)},x=1in,y=-1in]
\node[anchor=north west,inner sep=0pt,text width=7.2in,align=left,font=\fontsize{11}{13.75}\selectfont] at (0.65,0.41) {Week 19 / No card inside another / K--1};
\draw[line width=.35pt] (.65,.70)--(7.85,.70);
\node[anchor=north west,inner sep=0pt,text width=6.7in,align=left,font=\fontsize{9}{11.25}\selectfont] at (0.65,10.43) {Bellingham Math Circle / Week 19 / F19-K-v2};
\node[anchor=north west,inner sep=0pt,text width=0.55in,align=right,font=\fontsize{9}{11.25}\selectfont] at (7.3,10.43) {3};
\node[anchor=north west,inner sep=0pt,text width=7.2in,align=left,font=\fontsize{16}{20.0}\selectfont] at (0.65,0.98) {\textbf{Problem 4:} Make the biggest collection you can from this deck. Make a different collection just as big.};
\draw[line width=.7pt,rounded corners=2pt] (0.79,1.98) rectangle (1.48,2.67);
\draw[line width=1.5pt] (0.8452000000000001,1.98) -- (0.9004000000000001,1.98);
\draw[line width=.95pt] (0.9832000000000001,2.1732) circle[radius=0.08417999999999999in];
\draw[line width=.95pt] (0.89902,2.39262) rectangle (1.06738,2.56098);
\draw[line width=.7pt,rounded corners=2pt] (1.68,1.98) rectangle (2.37,2.67);
\draw[line width=1.5pt] (1.7351999999999999,1.98) -- (1.7904,1.98);
\draw[line width=.7pt,rounded corners=2pt] (2.57,1.98) rectangle (3.26,2.67);
\draw[line width=1.5pt] (2.6252,1.98) -- (2.6803999999999997,1.98);
\draw[line width=.95pt,line join=round] (3.0667999999999997,2.076393)--(3.159398,2.244753)--(2.9742019999999996,2.244753)--cycle;
\draw[line width=.7pt,rounded corners=2pt] (3.46,1.98) rectangle (4.15,2.67);
\draw[line width=1.5pt] (3.5152,1.98) -- (3.5704,1.98);
\draw[line width=.95pt] (3.6532,2.1732) circle[radius=0.08417999999999999in];
\draw[line width=.95pt,line join=round] (3.9568,2.076393)--(4.049398,2.244753)--(3.8642019999999997,2.244753)--cycle;
\draw[line width=.95pt] (3.56902,2.39262) rectangle (3.73738,2.56098);
\draw[line width=.7pt,rounded corners=2pt] (4.35,1.98) rectangle (5.039999999999999,2.67);
\draw[line width=1.5pt] (4.4052,1.98) -- (4.4604,1.98);
\draw[line width=.95pt] (4.5432,2.1732) circle[radius=0.08417999999999999in];
\draw[line width=.7pt,rounded corners=2pt] (5.239999999999999,1.98) rectangle (5.93,2.67);
\draw[line width=1.5pt] (5.2951999999999995,1.98) -- (5.3504,1.98);
\draw[line width=.95pt,line join=round] (5.7368,2.076393)--(5.829397999999999,2.244753)--(5.644202,2.244753)--cycle;
\draw[line width=.95pt] (5.349019999999999,2.39262) rectangle (5.517379999999999,2.56098);
\draw[line width=.7pt,rounded corners=2pt] (6.13,1.98) rectangle (6.82,2.67);
\draw[line width=1.5pt] (6.1852,1.98) -- (6.2404,1.98);
\draw[line width=.95pt] (6.3232,2.1732) circle[radius=0.08417999999999999in];
\draw[line width=.95pt,line join=round] (6.6268,2.076393)--(6.719398,2.244753)--(6.5342020000000005,2.244753)--cycle;
\draw[line width=.7pt,rounded corners=2pt] (7.02,1.98) rectangle (7.709999999999999,2.67);
\draw[line width=1.5pt] (7.0752,1.98) -- (7.1304,1.98);
\draw[line width=.95pt] (7.12902,2.39262) rectangle (7.2973799999999995,2.56098);
\draw[black!22,line width=.35pt] (.65,4.77)--(7.85,4.77);
\node[anchor=north west,inner sep=0pt,text width=7.2in,align=left,font=\fontsize{16}{20.0}\selectfont] at (0.65,5.04) {\textbf{Problem 5:} Use the same deck. Keep each group of cards below and add as many cards as you can.};
\draw[line width=.7pt,rounded corners=2pt] (0.97,6.1) rectangle (1.75,6.88);
\draw[line width=1.5pt] (1.0324,6.1) -- (1.0948,6.1);
\draw[line width=.95pt] (1.1884000000000001,6.3184) circle[radius=0.09516in];
\draw[line width=.95pt,line join=round] (1.5316,6.208965999999999)--(1.636276,6.399285999999999)--(1.426924,6.399285999999999)--cycle;
\draw[line width=.95pt] (1.0932400000000002,6.56644) rectangle (1.28356,6.75676);
\draw[line width=.7pt,rounded corners=2pt] (4.7,6.1) rectangle (5.48,6.88);
\draw[line width=1.5pt] (4.7624,6.1) -- (4.8248,6.1);
\draw[line width=.95pt] (4.9184,6.3184) circle[radius=0.09516in];
\draw[line width=.7pt,rounded corners=2pt] (0.97,8.05) rectangle (1.75,8.83);
\draw[line width=1.5pt] (1.0324,8.05) -- (1.0948,8.05);
\draw[line width=.95pt] (1.1884000000000001,8.268400000000002) circle[radius=0.09516in];
\draw[line width=.95pt,line join=round] (1.5316,8.158966000000001)--(1.636276,8.349286000000001)--(1.426924,8.349286000000001)--cycle;
\draw[line width=.7pt,rounded corners=2pt] (4.7,8.05) rectangle (5.48,8.83);
\draw[line width=1.5pt] (4.7624,8.05) -- (4.8248,8.05);
\draw[line width=.95pt] (4.9184,8.268400000000002) circle[radius=0.09516in];
\draw[line width=.7pt,rounded corners=2pt] (5.7,8.05) rectangle (6.48,8.83);
\draw[line width=1.5pt] (5.7624,8.05) -- (5.8248,8.05);
\draw[line width=.95pt,line join=round] (6.2616000000000005,8.158966000000001)--(6.366276000000001,8.349286000000001)--(6.156924,8.349286000000001)--cycle;
\draw[line width=.95pt] (5.82324,8.516440000000001) rectangle (6.01356,8.706760000000001);
\end{scope}\end{tikzpicture}\mbox{}
\newpage
\begin{tikzpicture}[remember picture,overlay]\begin{scope}[shift={(current page.north west)},x=1in,y=-1in]
\node[anchor=north west,inner sep=0pt,text width=7.2in,align=left,font=\fontsize{11}{13.75}\selectfont] at (0.65,0.41) {Week 19 / No card inside another / K--1};
\draw[line width=.35pt] (.65,.70)--(7.85,.70);
\node[anchor=north west,inner sep=0pt,text width=6.7in,align=left,font=\fontsize{9}{11.25}\selectfont] at (0.65,10.43) {Bellingham Math Circle / Week 19 / F19-K-v2};
\node[anchor=north west,inner sep=0pt,text width=0.55in,align=right,font=\fontsize{9}{11.25}\selectfont] at (7.3,10.43) {4};
\node[anchor=north west,inner sep=0pt,text width=7.2in,align=left,font=\fontsize{16}{20.0}\selectfont] at (0.65,0.98) {\textbf{Problem 6:} For each deck, use all the cards to make as few rows as you can. Each card fits inside the next; a card can be alone.};
\draw[line width=.7pt,rounded corners=2pt] (1.0,2.06) rectangle (1.62,2.68);
\draw[line width=1.5pt] (1.0496,2.06) -- (1.0992,2.06);
\draw[line width=.95pt] (1.1736,2.2336) circle[radius=0.07564in];
\draw[line width=.7pt,rounded corners=2pt] (2.43,2.06) rectangle (3.0500000000000003,2.68);
\draw[line width=1.5pt] (2.4796,2.06) -- (2.5292000000000003,2.06);
\draw[line width=.95pt] (2.6036,2.2336) circle[radius=0.07564in];
\draw[line width=.95pt,line join=round] (2.8764000000000003,2.146614)--(2.959604,2.297894)--(2.7931960000000005,2.297894)--cycle;
\draw[-stealth,line width=.8pt] (1.74,2.37)--(2.30,2.37);
\node[anchor=north west,inner sep=0pt,text width=4.1in,align=left,font=\fontsize{12}{15.0}\selectfont] at (3.4,2.02) {Example of one nested row: the first card fits inside the next. This is not a whole-deck arrangement.};
\draw[line width=.7pt,rounded corners=2pt] (2.32,3.17) rectangle (3.09,3.94);
\draw[line width=1.5pt] (2.3815999999999997,3.17) -- (2.4432,3.17);
\draw[line width=.7pt,rounded corners=2pt] (3.3499999999999996,3.17) rectangle (4.119999999999999,3.94);
\draw[line width=1.5pt] (3.4115999999999995,3.17) -- (3.4732,3.17);
\draw[line width=.95pt] (3.5656,3.3856) circle[radius=0.09394in];
\draw[line width=.95pt,line join=round] (3.9044,3.277569)--(4.007734,3.465449)--(3.801066,3.465449)--cycle;
\draw[line width=.7pt,rounded corners=2pt] (4.38,3.17) rectangle (5.15,3.94);
\draw[line width=1.5pt] (4.4416,3.17) -- (4.5032,3.17);
\draw[line width=.95pt] (4.5956,3.3856) circle[radius=0.09394in];
\draw[line width=.7pt,rounded corners=2pt] (5.41,3.17) rectangle (6.18,3.94);
\draw[line width=1.5pt] (5.4716000000000005,3.17) -- (5.5332,3.17);
\draw[line width=.95pt,line join=round] (5.9644,3.277569)--(6.067734000000001,3.465449)--(5.861066,3.465449)--cycle;
\draw[black!22,line width=.35pt] (.65,5.24)--(7.85,5.24);
\draw[line width=.7pt,rounded corners=2pt] (0.79,5.58) rectangle (1.48,6.27);
\draw[line width=1.5pt] (0.8452000000000001,5.58) -- (0.9004000000000001,5.58);
\draw[line width=.95pt] (0.9832000000000001,5.7732) circle[radius=0.08417999999999999in];
\draw[line width=.95pt] (0.89902,5.9926200000000005) rectangle (1.06738,6.16098);
\draw[line width=.7pt,rounded corners=2pt] (1.68,5.58) rectangle (2.37,6.27);
\draw[line width=1.5pt] (1.7351999999999999,5.58) -- (1.7904,5.58);
\draw[line width=.7pt,rounded corners=2pt] (2.57,5.58) rectangle (3.26,6.27);
\draw[line width=1.5pt] (2.6252,5.58) -- (2.6803999999999997,5.58);
\draw[line width=.95pt,line join=round] (3.0667999999999997,5.676393)--(3.159398,5.844753)--(2.9742019999999996,5.844753)--cycle;
\draw[line width=.7pt,rounded corners=2pt] (3.46,5.58) rectangle (4.15,6.27);
\draw[line width=1.5pt] (3.5152,5.58) -- (3.5704,5.58);
\draw[line width=.95pt] (3.6532,5.7732) circle[radius=0.08417999999999999in];
\draw[line width=.95pt,line join=round] (3.9568,5.676393)--(4.049398,5.844753)--(3.8642019999999997,5.844753)--cycle;
\draw[line width=.95pt] (3.56902,5.9926200000000005) rectangle (3.73738,6.16098);
\draw[line width=.7pt,rounded corners=2pt] (4.35,5.58) rectangle (5.039999999999999,6.27);
\draw[line width=1.5pt] (4.4052,5.58) -- (4.4604,5.58);
\draw[line width=.95pt] (4.5432,5.7732) circle[radius=0.08417999999999999in];
\draw[line width=.7pt,rounded corners=2pt] (5.239999999999999,5.58) rectangle (5.93,6.27);
\draw[line width=1.5pt] (5.2951999999999995,5.58) -- (5.3504,5.58);
\draw[line width=.95pt,line join=round] (5.7368,5.676393)--(5.829397999999999,5.844753)--(5.644202,5.844753)--cycle;
\draw[line width=.95pt] (5.349019999999999,5.9926200000000005) rectangle (5.517379999999999,6.16098);
\draw[line width=.7pt,rounded corners=2pt] (6.13,5.58) rectangle (6.82,6.27);
\draw[line width=1.5pt] (6.1852,5.58) -- (6.2404,5.58);
\draw[line width=.95pt] (6.3232,5.7732) circle[radius=0.08417999999999999in];
\draw[line width=.95pt,line join=round] (6.6268,5.676393)--(6.719398,5.844753)--(6.5342020000000005,5.844753)--cycle;
\draw[line width=.7pt,rounded corners=2pt] (7.02,5.58) rectangle (7.709999999999999,6.27);
\draw[line width=1.5pt] (7.0752,5.58) -- (7.1304,5.58);
\draw[line width=.95pt] (7.12902,5.9926200000000005) rectangle (7.2973799999999995,6.16098);
\end{scope}\end{tikzpicture}\mbox{}
\newpage
\begin{tikzpicture}[remember picture,overlay]\begin{scope}[shift={(current page.north west)},x=1in,y=-1in]
\node[anchor=north west,inner sep=0pt,text width=7.2in,align=left,font=\fontsize{11}{13.75}\selectfont] at (0.65,0.41) {Week 19 / No card inside another / K--1};
\draw[line width=.35pt] (.65,.70)--(7.85,.70);
\node[anchor=north west,inner sep=0pt,text width=6.7in,align=left,font=\fontsize{9}{11.25}\selectfont] at (0.65,10.43) {Bellingham Math Circle / Week 19 / F19-K-v2};
\node[anchor=north west,inner sep=0pt,text width=0.55in,align=right,font=\fontsize{9}{11.25}\selectfont] at (7.3,10.43) {5};
\node[anchor=north west,inner sep=0pt,text width=7.2in,align=left,font=\fontsize{16}{20.0}\selectfont] at (0.65,0.98) {\textbf{Problem 7:} Can you choose four cards from this deck without breaking the rule? Show why or why not.};
\draw[line width=.7pt,rounded corners=2pt] (0.79,2.01) rectangle (1.48,2.6999999999999997);
\draw[line width=1.5pt] (0.8452000000000001,2.01) -- (0.9004000000000001,2.01);
\draw[line width=.95pt] (0.9832000000000001,2.2032) circle[radius=0.08417999999999999in];
\draw[line width=.95pt] (0.89902,2.4226199999999998) rectangle (1.06738,2.5909799999999996);
\draw[line width=.7pt,rounded corners=2pt] (1.68,2.01) rectangle (2.37,2.6999999999999997);
\draw[line width=1.5pt] (1.7351999999999999,2.01) -- (1.7904,2.01);
\draw[line width=.7pt,rounded corners=2pt] (2.57,2.01) rectangle (3.26,2.6999999999999997);
\draw[line width=1.5pt] (2.6252,2.01) -- (2.6803999999999997,2.01);
\draw[line width=.95pt,line join=round] (3.0667999999999997,2.1063929999999997)--(3.159398,2.274753)--(2.9742019999999996,2.274753)--cycle;
\draw[line width=.7pt,rounded corners=2pt] (3.46,2.01) rectangle (4.15,2.6999999999999997);
\draw[line width=1.5pt] (3.5152,2.01) -- (3.5704,2.01);
\draw[line width=.95pt] (3.6532,2.2032) circle[radius=0.08417999999999999in];
\draw[line width=.95pt,line join=round] (3.9568,2.1063929999999997)--(4.049398,2.274753)--(3.8642019999999997,2.274753)--cycle;
\draw[line width=.95pt] (3.56902,2.4226199999999998) rectangle (3.73738,2.5909799999999996);
\draw[line width=.7pt,rounded corners=2pt] (4.35,2.01) rectangle (5.039999999999999,2.6999999999999997);
\draw[line width=1.5pt] (4.4052,2.01) -- (4.4604,2.01);
\draw[line width=.95pt] (4.5432,2.2032) circle[radius=0.08417999999999999in];
\draw[line width=.7pt,rounded corners=2pt] (5.239999999999999,2.01) rectangle (5.93,2.6999999999999997);
\draw[line width=1.5pt] (5.2951999999999995,2.01) -- (5.3504,2.01);
\draw[line width=.95pt,line join=round] (5.7368,2.1063929999999997)--(5.829397999999999,2.274753)--(5.644202,2.274753)--cycle;
\draw[line width=.95pt] (5.349019999999999,2.4226199999999998) rectangle (5.517379999999999,2.5909799999999996);
\draw[line width=.7pt,rounded corners=2pt] (6.13,2.01) rectangle (6.82,2.6999999999999997);
\draw[line width=1.5pt] (6.1852,2.01) -- (6.2404,2.01);
\draw[line width=.95pt] (6.3232,2.2032) circle[radius=0.08417999999999999in];
\draw[line width=.95pt,line join=round] (6.6268,2.1063929999999997)--(6.719398,2.274753)--(6.5342020000000005,2.274753)--cycle;
\draw[line width=.7pt,rounded corners=2pt] (7.02,2.01) rectangle (7.709999999999999,2.6999999999999997);
\draw[line width=1.5pt] (7.0752,2.01) -- (7.1304,2.01);
\draw[line width=.95pt] (7.12902,2.4226199999999998) rectangle (7.2973799999999995,2.5909799999999996);
\draw[black!22,line width=.35pt] (.65,5.54)--(7.85,5.54);
\node[anchor=north west,inner sep=0pt,text width=7.2in,align=left,font=\fontsize{16}{20.0}\selectfont] at (0.65,5.82) {\textbf{Problem 8:} Set aside any one card from this deck. Can you still make a collection of three cards, whichever card you set aside?};
\end{scope}\end{tikzpicture}\mbox{}
\end{document}
